\chapter{What it means to complete a local CFT}
\label{ch:sewing}

A modular-invariant partition function is a powerful constraint,
but it is still a count of states. A local theory must also assign
correlation functions consistently when surfaces are cut and sewn.
This chapter explains that additional structure through a completely
worked Ising example, then returns to the Haagerup construction and
the reduced-theory proposals.

\section{Cutting a surface introduces a channel basis}
\index{sewing}\index{conformal block!pairing}

Cut a four-punctured sphere along a circle separating two insertions
from the other two. Insert a complete set of intermediate states
on that circle. In a rational chiral theory this produces a finite
sum of conformal blocks, each block already summing infinitely
many descendants in an intermediate module. If the relevant
intertwiner space has dimension greater than one, the channel
also carries a multiplicity index.

Let $\mathcal F_p^{(s)}(z)$ be a chosen basis of holomorphic
blocks for one cut. A full correlator takes the form
\begin{equation}
 G(z,\bar z)=\sum_{p,q}D_{pq}\,
       \mathcal F_p^{(s)}(z)\overline{\mathcal F_q^{(s)}(z)}.
 \label{eq:sewing-block-pairing}
\end{equation}
Here $p,q$ include multiplicity labels, and $D$ contains products
of normalized bulk three-point coefficients. It is not fixed by
the torus matrix $Z_{ij}$ alone.

Changing the cut changes the block basis by a fusing matrix.
Analytically continuing an insertion around another changes it by
a braiding or monodromy matrix. The pairing $D$ must transform
consistently so that the full correlator is single-valued and
independent of the chosen cut. In a reflection-positive channel,
the two groups of insertions are related by reflection and adjoint.
If one group prepares states $\Psi_a$, the reflected pairing is a
Gram matrix:
\[
 K_{ab}=\langle\Psi_a,\Psi_b\rangle,
 \qquad
 \sum_{a,b}\overline{c_a}K_{ab}c_b
 =\Bigl\|\sum_b c_b\Psi_b\Bigr\|^2\ge0.
\]
This explains the positive semidefinite coefficient pairing in
such a channel, with separated insertions and positive Euclidean
propagation understood. A basis change can move phases between
blocks and coefficients, but cannot remove a failure of these
compatibility conditions.

The categorical associator and braiding record these changes of
chiral channel. They constrain a local completion, while the
choice of bulk algebra and positive pairing specifies how the
left and right chiral data are combined.
For $U$ conformal blocks, these fusing matrices belong to
$\Rep(U)\simeq Z(\mathcal H_3)$. The associators for junctions
of the six selected symmetry lines belong to $\mathcal H_3$.
The center construction relates the two categories, but their
$F$ matrices are different objects. In particular the four-by-four
$\rho$ recoupling block is not a conformal-block matrix for
arbitrary $U$ primaries.

\section{A complete two-channel example from the Ising model}
\label{sec:sewing-Ising}
\index{Ising model}\index{BPZ equation}

The Ising spin field $\sigma$ has chiral weight $1/16$, and
$\sigma\times\sigma=\one+\epsilon$, with
$h_\epsilon=1/2$. We normalize the full fields by
\[
 \langle\sigma(z,\bar z)\sigma(0)\rangle=|z|^{-1/4},
 \qquad
 \langle\epsilon(z,\bar z)\epsilon(0)\rangle=|z|^{-2}.
\]
The full spin and energy fields have weights $(1/16,1/16)$
and $(1/2,1/2)$ respectively. We calculate
$G_\sigma=\langle\sigma(\infty)\sigma(1)\sigma(z,\bar z)
\sigma(0)\rangle$, where
$\sigma(\infty)=\lim_{|w|\to\infty}|w|^{1/4}\sigma(w,\bar w)$.
A useful pair of chiral blocks is
\begin{align}
 \mathcal F_1(z)&=\frac{[z(1-z)]^{-1/8}}{\sqrt2}
                         \sqrt{1+\sqrt{1-z}},\notag\\
 \mathcal F_\epsilon(z)&=\frac{[z(1-z)]^{-1/8}}{\sqrt2}
                         \sqrt{1-\sqrt{1-z}}.
 \label{eq:sewing-Ising-blocks}
\end{align}
Choose positive square roots on $0<z<1$ and then specify paths
for analytic continuation. Near zero,
\[
 \mathcal F_1(z)=z^{-1/8}(1+O(z^2)),\qquad
 \mathcal F_\epsilon(z)=\frac12z^{3/8}(1+O(z)).
\]
The exponent difference is the exchanged weight $1/2$.

These blocks can be obtained from the level-two null vector of
$\sigma$. The Gram-matrix exercise gave
$(L_{-2}-\tfrac43L_{-1}^2)|\sigma\rangle=0$.
Insert this null field in a correlator and move the stress contour
onto the other insertions. The Ward identity replaces the
$L_{-2}$ insertion by
$\sum_{i\ne z}[h_\sigma/(z-z_i)^2+
(z-z_i)^{-1}\partial_{z_i}]$ acting on the correlator.
Use the three global conformal Ward identities to fix the other
insertions at $0,1,\infty$. Their derivatives become
$\partial_0 f=(z-1)f'+2h_\sigma f$ and
$\partial_1 f=-zf'-2h_\sigma f$ before substituting the
fixed positions. For
$f(z)=\langle\sigma(\infty)\sigma(1)\sigma(z)\sigma(0)\rangle$
the resulting differential equation is
\begin{equation}
 f''(z)+\frac{3(1-2z)}{4z(1-z)}f'(z)
       -\frac{3}{64z^2(1-z)^2}f(z)=0.
 \label{eq:sewing-Ising-BPZ}
\end{equation}
Its indicial equation at zero is
$r(r-1)+3r/4-3/64=0$, with roots $-1/8$ and $3/8$.
They are precisely $h_{\rm exchange}-2h_\sigma$ for the
vacuum and energy channels. Substitution verifies the two
radical solutions above. A second-order equation has a
two-dimensional local solution space away from its singularities,
so these independent solutions exhaust the chiral blocks for
this null-vector problem.

Under exchange of the $s$ and $t$ cuts,
\begin{equation}
 \begin{pmatrix}\mathcal F_1(1-z)\\\mathcal F_\epsilon(1-z)\end{pmatrix}
 =\frac1{\sqrt2}
  \begin{pmatrix}1&1\\1&-1\end{pmatrix}
  \begin{pmatrix}\mathcal F_1(z)\\\mathcal F_\epsilon(z)\end{pmatrix}.
 \label{eq:sewing-Ising-F}
\end{equation}
The identity follows by squaring the half-angle radical relations
and checking their signs on $0<z<1$. The matrix is real orthogonal.
It does not preserve either individual channel: a vacuum block
in one cut contains both channels in the other.

Suppose the full correlator has pairing matrix $D$. Monodromy
around zero gives phases $e^{-\pi i/4}$ and $e^{3\pi i/4}$ to
the two holomorphic blocks. The phases differ by a sign, so
single-valuedness removes off-diagonal entries of $D$.
Write $D=\diag(a,b)$. Crossing requires
$F^{\mathsf t}DF=D$, and its off-diagonal entry is
$(a-b)/2$. Thus $a=b$. Unit normalization of the vacuum channel
sets $a=1$, yielding
\begin{equation}
 G_\sigma(z,\bar z)=|\mathcal F_1(z)|^2+
                              |\mathcal F_\epsilon(z)|^2.
 \label{eq:sewing-Ising-full}
\end{equation}
This is a local, crossing-symmetric positive pairing. In the
$z\to0$ expansion, the energy channel has relative coefficient
$1/4$, so the normalized full coefficient obeys
$|C_{\sigma\sigma\epsilon}|=1/2$.

The final answer also has a branch-independent form:
\begin{equation}
 G_\sigma(z,\bar z)=
 \frac{\sqrt{1+|z|+|1-z|}}
      {\sqrt2\,|z(1-z)|^{1/4}}.
 \label{eq:sewing-Ising-singlevalued}
\end{equation}
Indeed, for $w=\sqrt{1-z}$ one has
$(|1+w|+|1-w|)^2=2(1+|1-z|+|z|)$.
The expression makes single-valuedness explicit and verifies
both crossing generators:
\[
 G_\sigma(1-z,1-\bar z)=G_\sigma(z,\bar z),\qquad
 G_\sigma(z,\bar z)=|z|^{-1/4}
             G_\sigma(1/z,1/\bar z).
\]
The prefactor in the second equation is fixed by the external
dimension; crossing is not invariance under every fractional
linear transformation without such factors.

Channel mixing explains why a single chiral block usually does not
give a full correlator by taking its absolute square. Here
$|\mathcal F_1|^2$ alone fails crossing, even though it is
nonnegative on the real interval and has the correct leading
vacuum singularity. Those two properties are insufficient.

\section{Algebra objects organize rational sewing}
\index{Frobenius algebra}\index{full center}\index{Morita equivalence}

In the rational framework, a suitable symmetric special Frobenius
algebra object $A$ in the representation category of an existing
rational chiral theory organizes the oriented full-CFT sewing
data. The chiral theory and its analytic conformal blocks are
inputs; an abstract modular category alone does not construct
them. Multiplication specifies how algebra lines meet,
associativity makes different decompositions agree, and the
Frobenius relation relates multiplication to a compatible
comultiplication. In the unitary Q-system normalization this
comultiplication is the adjoint of multiplication.
The positive Q-system and unitary chiral theory provide the adjoint
and Hilbert-space pairings; the purely algebraic Frobenius identities
alone do not assert reflection positivity. The resulting bulk state
space is described by the full center of $A$.

The TFT construction proves mapping-class invariance and bulk and
boundary factorization for these categorical correlators on oriented
world sheets. Their interpretation as analytic amplitudes on Riemann
surfaces also uses the chiral conformal-block theory, including its
sewing and convergence theorems. The annulus comparison in
Chapter~\ref{ch:unitarity} establishes the positivity and VOA--net
statements used here; it is not itself a proof of every higher-genus
analytic sewing theorem. The diagonal local conformal net already
exists by the construction in Chapter~\ref{ch:full_theories}.

For the Cardy or diagonal construction, $A=\one$. Its bulk
object is the canonical sum
\[
 \bigoplus_i M_i\boxtimes M_i^*
\]
in $\Rep(U)\boxtimes\Rep(U)^{\rm rev}$, with a multiplication
built from the chiral fusion and duality data. The torus trace is
the diagonal partition function, with the right-moving
contragredient convention understood. All twelve Haagerup labels,
and all Ising labels used here, are self-dual, so the distinction
between diagonal and charge-conjugate labelings does not change
these examples. For a local chiral extension,
a connected commutative separable twist-free algebra in $\Rep(U)$,
with the positive structure and grading conditions of the extension
theorem, supplies additional chiral fields; the selected Haagerup
Lagrangian algebra gives the
holomorphic extension studied above.

The underlying object is only part of the algebra. A decomposition
such as $A=\one\oplus b\oplus2a$ does not state its multiplication
maps or their associativity. Equal characters of $a$ and $c_0$
make this particularly visible: a scalar trace can conceal an
extension choice that the junction data distinguish.

Morita-equivalent algebra objects give equivalent full-CFT
descriptions in the appropriate rational construction, while
their boundary presentations can look different. This is one
reason module categories, rather than only fusion rings, are
natural physical objects. The associator determines whether a
candidate multiplication is actually associative; a list of
integer fusion multiplicities cannot replace it.

\section{Existence and explicit OPE tables are different achievements}

Existence of the diagonal full theory does not require printing
every coefficient in a basis of its bulk primary fields.
Conversely, a long table of low
three-point coefficients does not by itself prove all sewing
identities in all sectors and genera. These are complementary
forms of information.

For the Haagerup $c=8$ construction, the lattice vacuum algebra
supplies explicit chiral OPEs, the complete weight-two primary
tensor, and several exact four-point functions. The structural
and analytic arguments supply the representation category,
unitarity, and full-theory existence under their stated inputs.
The selected lattice intertwiners in the lightest module give
additional nonvacuum chiral data.

What is not supplied by those particular calculations is a
complete basis of all nonvacuum intertwiner spaces together with
the numerical or algebraic bulk sewing tensor for every diagonal
sector. In particular the nine simple lattice blocks of
\eqref{eq:corr-light-block} are not a declaration of the complete
lightest-scalar bulk four-point function. They remember the
containing $B$ extension and its chosen intertwiners.

For full vacuum-sector fields that factor as
$\Phi(z,\bar z)=v(z)\bar w(\bar z)$, correlators do factor into
the corresponding chiral and antiholomorphic correlators. This
gives many explicit local observables in both full theories.
The factorization follows from the specified tensor-product
sector, not from a general rule that every nonvacuum bulk field
is the product of any two selected chiral blocks.

\section{A coset is a chiral operation}
\index{coset!full-theory completion}

Suppose $A$ and $C=\Com(A,V)$ commute inside a VOA $V$,
with complementary conformal vectors
$\omega_V=\omega_A+\omega_C$.
When an ambient module has a semisimple restriction with
compatible ordinary gradings, it can branch as
\begin{equation}
 V_i\downarrow_{A\otimes C}
 =\bigoplus_\lambda A_\lambda\otimes B_{i\lambda},
 \qquad
 \chi_i^V=\sum_\lambda\chi_\lambda^A b_{i\lambda}.
 \label{eq:sewing-branching}
\end{equation}
Here $B_{i\lambda}$ is a multiplicity space carrying $C$ fields,
and $b_{i\lambda}$ is its graded trace when that trace is defined.
The equality of characters is a branching identity. Deleting the
$A$ characters from a full partition function is not a general
sewing theorem.

Several issues can arise. A multiplicity space might fail to be
irreducible or ordinary for $C$. Infinitely many branching labels
can produce a divergent reduced thermal trace. Distinct ambient
sectors can produce repeated copies of the reduced vacuum. A
putative left/right pairing can fail modular invariance or crossing.
Finally, an ambient defect need not descend with coherent junctions.
These are separate questions and can fail independently.

The exact reverse-commutant example in
Chapter~\ref{ch:cosets} gives a memorable finite warning: a naive
branching sum produces $6Z_L$, because multiplicities contribute
$1^2+1^2+2^2=6$ vacuum copies. The lattice CFT $Z_L$ may exist,
but obtaining it requires a new, justified choice of completion.
Dividing a trace by six does not construct its multiplication or
its defect action.

For an irrational stress splitting, the challenge can be analytic
as well as categorical. A formal diagonal pairing of branching
spaces with one vacuum may still have infinitely many states
below some energy and an infinite thermal trace. Conversely,
projecting to a single chiral vacuum character can violate modular
invariance. A continuous modular density with no normalizable
vacuum describes yet another possibility and has different
effective-central-charge behavior.

\section{Defect descent requires junctions}
\index{defect!descent}\index{junction}

A line preserving a candidate stress tensor preserves its local
conformal transformations. This is necessary for a topological
symmetry of the reduced theory. It is not sufficient to identify
the reduced defect category with $\mathcal H_3$.
One needs images of the lines and of their junction spaces,
compatible tensor-product maps, and the same coherent associativity
relations. In categorical language, one needs an appropriate
monoidal transport, not only a list of operators commuting with
$L_0$.

For example, operators may satisfy the fusion-algebra relations,
including $\rho^2=1+\rho+\alpha\rho+\alpha^2\rho$, on a selected
vector space while carrying no information about the $F$ symbols of
four line junctions. The ordinary Haagerup fusion algebra has
only two one-dimensional representations and one two-dimensional
representation. It cannot encode the full tensor product of
local charge sectors by itself. The tube algebra and its
half-braiding data supply more information, but their transport
must still be tied to the actual reduced fields and sewing.

This distinction is constructive rather than merely restrictive.
In the rational $c=6$ commutant discussed later, the extension
correspondence transports an algebra object and its multiplication.
That supplies junction data, which is stronger than a numerical
agreement of line eigenvalues. For the irrational near-two axis,
an analogous full transport has not been completed in the
research record.

\section{How a proposed identification should be assessed}

A convincing identification of a microscopic continuum theory
should specify the actual Hilbert space and vacuum, a convergent
torus trace, normalized local fields and their OPEs, crossing and
sewing, and the intended symmetry action including junctions.
One rarely computes all of these from a lattice simulation.
Instead, an exact construction can supply them and numerical
data can test whether the lattice approaches that construction.

Matching central charge is the first coarse test. Matching
dimensions, spins, multiplicities, and symmetry sectors is much
stronger. Matching normalized OPE ratios and phase invariants is
stronger again. Every comparison needs control of finite-size,
finite-entanglement, and operator-normalization errors. A mismatch
can exclude a specified candidate under explicit limiting
assumptions, but an unexplained mismatch need not classify all
possible CFTs at the same central charge.

This is the perspective we will use for the physical Haagerup
chain. It makes the existing evidence informative without turning
a suggestive decimal or a chiral stress vector into a completed
bulk construction.

\section{Exercises}

\begin{exercise}\label{ex:sewing-Ising}
Verify \eqref{eq:sewing-Ising-F} on the interval $0<z<1$ by
elementary radical identities. Use it to show directly that
$|\mathcal F_1|^2$ alone fails crossing.

\par\smallskip\noindent{\footnotesize\normalfont
\hyperref[sol:ex:sewing-Ising]{Solution on page~\pageref*{sol:ex:sewing-Ising}.}}\par
\end{exercise}

\begin{exercise}\label{ex:sewing-pairing}
Starting with a general Hermitian $2\times2$ matrix $D$, use
the two Ising monodromy phases and the fusing matrix to determine
all local crossing-invariant pairings. Impose positivity and
vacuum normalization at the end.

\par\smallskip\noindent{\footnotesize\normalfont
\hyperref[sol:ex:sewing-pairing]{Solution on page~\pageref*{sol:ex:sewing-pairing}.}}\par
\end{exercise}

\begin{exercise}\label{ex:sewing-object}
Explain the difference between an equality of characters, an
isomorphism of underlying modules, an isomorphism of algebra
objects, and an equivalence of full local CFTs. Give one piece
of data forgotten at each step toward the character.

\par\smallskip\noindent{\footnotesize\normalfont
\hyperref[sol:ex:sewing-object]{Solution on page~\pageref*{sol:ex:sewing-object}.}}\par
\end{exercise}

\begin{exercise}\label{ex:sewing-coset}
A proposed reduced theory has one vacuum but infinitely many
states with energies bounded by a fixed constant. Prove that
its thermal trace diverges at every positive inverse temperature.
Why does positivity make this argument especially direct?

\par\smallskip\noindent{\footnotesize\normalfont
\hyperref[sol:ex:sewing-coset]{Solution on page~\pageref*{sol:ex:sewing-coset}.}}\par
\end{exercise}

\begin{exercise}\label{ex:sewing-descent}
An ambient Haagerup line action commutes with two complementary
stress tensors. List the extra structures needed to turn its
restriction into a specified Haagerup symmetry of one reduced
full theory. Which of these are visible in a torus trace?

\par\smallskip\noindent{\footnotesize\normalfont
\hyperref[sol:ex:sewing-descent]{Solution on page~\pageref*{sol:ex:sewing-descent}.}}\par
\end{exercise}
